% Строчки для суммарной таблицы, в следующей строке прописаны используемые экземпляры узлов
%===s control, LLN0, GENRBRF1

%\needspace{4\baselineskip}
\color{uniblue}{\subsection{Устройство резервирования при отказе выключателя стороны ВН (УРОВ ВН)}}
\color{black}

\begin{enumerate}[label=\arabic{section}.\arabic{subsection}.\arabic*, labelsep=4pt, leftmargin=0pt, itemindent=57pt]

\item
Функция УРОВ предназначена для ликвидации повреждений, сопровождающихся отказом выключателя стороны ВН трансформатора.

\item
Основные характеристики функционального блока <<УРОВ ВН>>:
\begin{itemize}
    \item время срабатывания ИО тока при повышении тока скачком от 0 до $2I_{\textsl{уст.}}$ не более 30 мс;
    \item время возврата ИО тока при снижении тока скачком от $25I_{\textsl{ном}}$ до 0 не более 25 мс.
\end{itemize}

\item\label{urov2:punkt}
В случае присутствия сигналов срабатывания защит от логики отключения стороны ВН трансформатора при введенном состоянии функция <<УРОВ ВН>>: 
\begin{itemize}
\item формирует сигнал <<УРОВ: Пуск>> в следующих случаях:
	\begin{itemize}
	\sloppy
	\item если программный переключатель SGF3 <<Подхват\_по\_току>> находится в состоянии <<Не предусмотрено>> и через выключатель протекает ток, превышающий заданную уставку срабатывания ИО максимального тока функции <<УРОВ>> (в этом случае при пропадании сигнала пуска от защит возврат <<УРОВ>> может произойти раньше, чем наберется выдержка времени таймера T1 (уставка <<Тср>>), хотя ток через выключатель продолжает протекать); 
	\item если программный переключатель SGF3 <<Подхват\_по\_току>> находится в состоянии <<Предусмотрено>> и через выключатель протекает ток, превышающий заданную уставку срабатывания ИО максимального тока функции <<УРОВ>> (в этом случае сигнал пуска от защит запоминается на триггере, который сбрасывается по факту пропадания тока через выключатель);
	\end{itemize}
\item формирует сигнал <<УРОВ: Срабатывание>> через выдержку таймера Т1 (уставка <<Тср>>) при выведенном программном переключателе SGF2 <<Уск\_при\_блк\_откл\_В>> или введенном программном переключателе SGF2 <<Уск\_при\_блк\_откл\_В>> и отсутствии сигнала <<КСВ: Блокировка откл>> на входе функции;
\item формирует сигналы <<УРОВ: Срабатывание>> и <<УРОВ: Ускорение>> без выдержки времени при введенном программном переключателе SGF2 <<Уск\_при\_блк\_откл\_В>> и активном состоянии сигнала <<КСВ: Блокировка откл>>;
\end{itemize}
\item
В случае присутствия сигналов срабатывания от внешних защит сторон НН трансформатора (активация сигнала <<Пуск УРОВ внешний>>) при введенном состоянии функция <<УРОВ ВН>>: 
\begin{itemize}
	\item формирует сигнал <<УРОВ: Сраб <<на себя>> для действия на отключение своего выключателя, при этом, в качестве дополнительного условия, есть возможность контроля протекания тока через выключатель от своих измерительных органов (программный переключатель SGF6 <<Контр\_тока\_на\_себя>> должен находиться в состоянии <<Предусмотрено по внутр. ПО>>) или от пусковых органов токовых защит в составе устройства (программный переключатель SGF6 <<Контр\_тока\_на\_себя>> должен находиться в состоянии <<Предусмотрено по внеш. ПО>>). В случае введенного состояния программного переключателя SGF2 <<Уск\_при\_блк\_откл\_В>> дополнительно контролируется отсутствие входного сигнала <<КСВ: Блокировка откл>>, т.к. в случае присутствия сигнала <<КСВ: Блокировка откл>> действие на отключение своего выключателя должно быть исключено;
	\item при введенном программном переключателе SGF5 <<Действ\_на\_выш\_выкл>> формирует сигналы <<УРОВ: Пуск>>, <<УРОВ: Срабатывание>> и <<УРОВ: Ускорение>> в случаях, описанных в пункте \ref{urov2:punkt}.
\end{itemize}
\item
В алгоритме <<УРОВ>> дополнительно предусмотрен контроль срабатывания электромагнитов отключения выключателя выбором режима программного переключателя SGF4 <<Контроль\_ЭМО>>:
\begin{itemize}
	\item при состоянии <<Не предусмотрено>> контроль ЭМО не происходит;
	\item при состоянии <<Предусмотрено по ЭМО1>> контролируется цепь электромагнита ЭМО1, в этом случае дополнительным разрешающим условием для работы УРОВ является неактивное состояние сигнала <<Контроль цепи ЭМО1>> (предполагается, что дискретный вход устройства для контроля цепи ЭМО1 шунтируется сработанным контактом выходного реле защит);
	\item при состоянии <<Предусмотрено по ЭМО1 и ЭМО2>> контролируются цепи электромагнитов ЭМО1 и ЭМО2, в этом случае дополнительным разрешающим условием для работы УРОВ является неактивное состояние любого из сигналов <<Контроль цепи ЭМО1>> или <<Контроль цепи ЭМО2>> (предполагается, что как минимум один из дискретных входов устройства для контроля цепей ЭМО шунтируется сработанным контактом выходного реле защит);
\end{itemize}

\item
Ввод функции в работу на этапе параметрирования устройства осуществляется программным переключателем SGF1 <<Ввод\_функции>>. Информация о введенной функции отображается выходным сигналом <<УРОВ: Ввод>>.
\item
Функция <<УРОВ ВН>> может быть выведена из работы оперативно путем активации сигнала <<ОВ УРОВ ВН>>, что характеризуется активным состоянием сигнала <<УРОВ: Оперативный вывод>> на выходе функции. 

\item
Алгоритм работы <<УРОВ>> выполнен в соответствии с рисунком \ref{urov35:img1}. В таблице \ref{urov35:tbl1} приведены параметры, необходимые для настройки функции <<УРОВ>>.

\vspace{3mm}
\begin{figure}[H]
\centering
\includegraphics[width=\textwidth,height=\textheight,keepaspectratio]{\fbpath/26. УРОВ 35/_latex/img/1.pdf}
\captionof{figure}{Функциональная схема <<УРОВ>>}\label{urov35:img1}
\end{figure}

\small
\begin{longtable}{|>{\centering\arraybackslash}m{5.4cm}|>{\centering\arraybackslash}m{3.6cm}|>{\centering\arraybackslash}m{4.2cm}|>{\centering\arraybackslash}m{1.8cm}|>{\centering\arraybackslash}m{1cm}|}
\caption{Параметры для настройки функции <<УРОВ>>\hfill\vspace{-0.5\baselineskip}}\label{urov35:tbl1}\\ 
\hline
\rowcolor{gray!30}
Параметр (Параметр на ИЧМ) & Условное обозначение на схеме & Значение/ Диапазон & Единица измерения & Шаг \\ 
\hline
\endfirsthead
%\caption{\emph{Продолжение\hfill\vspace{-0.5\baselineskip}}} \\ % Отступ обозначения таблицы от таблицы и выравнивание надписи слева
\caption*{\hspace{3pt}\emph{Продолжение таблицы \ref{urov35:tbl1}\hfill\vspace{-0.5\baselineskip}}} \\ % сделано по ГОСТ 2.105 п.6.8.7
\hline
\rowcolor{gray!30}
Параметр (Параметр на ИЧМ) & Условное обозначение на схеме & Значение/ Диапазон & Единица измерения & Шаг \\ 
%\hline
\endhead
%\hline
\endfoot
%\hline
\endlastfoot
%==+t1*GENRBRF1|T_TPBRF
%===t1*GENRBRF1
\centering Ввод функции в работу (Ввод\_функции) & \centering SGF1 & \centering 0 = Не предусмотрено\\1 = Предусмотрено & \centering -- & \centering \arraybackslash -- \\
\hline
\centering Ток срабатывания (Iср) & \centering I$>$ & \centering 0,05 ... 0,50 & \centering о.е. & \centering \arraybackslash 0,01 \\
\hline
\centering Ускорение при блокировке отключения В (Уск\_при\_блк\_откл\_В) & \centering SGF2 & \centering 0 = Не предусмотрено\\1 = Предусмотрено & \centering -- & \centering \arraybackslash -- \\
\hline
\centering УРОВ с подхватом по току (Подхват\_по\_току) & \centering SGF3 & \centering 0 = Не предусмотрено\\1 = Предусмотрено & \centering -- & \centering \arraybackslash -- \\
\hline
\centering Контроль по току при действии «на себя» (Контр\_тока\_на\_себя) & \centering SGF6 & \centering 0 = Не предусмотрено\\1 = Предусмотрено по внутр. ПО\\2 = Предусмотрено по внеш. ПО & \centering -- & \centering \arraybackslash -- \\
\hline
\centering Контроль ЭМО (Контроль\_ЭМО) & \centering SGF4 & \centering 0 = Не предусмотрено\\1 = Предусмотрено по ЭМО1\\2 = Предусмотрено по ЭМО1 и ЭМО2 & \centering -- & \centering \arraybackslash -- \\
\hline
\centering Действие внешнего УРОВ на вышестоящий выключатель (Действ\_на\_выш\_выкл) & \centering SGF5 & \centering 0 = Не предусмотрено\\1 = Предусмотрено & \centering -- & \centering \arraybackslash -- \\
\hline
\centering Выдержка времени срабатывания (Tср) & \centering T1 & \centering 0,00 ... 1,00 & \centering с & \centering \arraybackslash 0,01 \\
\hline
%===t1
\end{longtable}
\normalsize

\end{enumerate}